% Circuito de medición para el Experimento 1: Onda cuadrada con DSO
\begin{circuitikz}[scale=0.9, transform shape]
    % Generador de funciones
    \draw (0,0) node[ground] {}
        to[sV, l={$G_1$ (\qty{1}{\kilo\hertz})}] (0,2)
        to[short] (1.5,2);
    \draw (0,0) to[short] (1.5,0);

    % Conexión y sonda atenuadora 10X
    \draw (1.5,2) to[R, l={Sonda 10X}] (4.5,2)
        to[short, -o] (5.2,2) node[above, font=\scriptsize\bfseries] {CH1};
    \draw (1.5,0) to[short] (4.5,0)
        to[short, -o] (5.2,0) node[above, font=\scriptsize\bfseries] {Masa};

    % Conexiones al panel del instrumento
    \draw (5.2,2) to[short] (5.8,2);
    \draw (5.2,0) to[short] (5.8,0);
    \draw (5.2,0) node[ground] {};

    % Bloque osciloscopio DSO
    \draw[thick, rounded corners=3pt, fill=blue!3] (5.8,-0.8) rectangle (10.4,2.8);
    \node at (8.1,2.3) [font=\bfseries\small] {Osciloscopio DSO};
    \node at (8.1,1.55) [font=\scriptsize, align=center] {Acoplamiento CC / Sin límite BW};
    \node at (8.1,0.95) [font=\scriptsize, align=center] {Math: FFT / Ventana Flattop};
    \node at (8.1,0.35) [font=\scriptsize, align=center] {Adquisición: Promedio 64};

    % Cuadro TP Calibración
    \node[draw, dashed, rounded corners, font=\tiny, align=center, fill=gray!8] at (8.1,-0.35) 
        {TP Calibración (\qty{1}{\kilo\hertz}, \qty{5}{\volt_{pp}})};
\end{circuitikz}
